\documentclass[11pt]{article}

\usepackage[T1]{fontenc}
\usepackage[utf8]{inputenc}
\usepackage{lmodern}
\usepackage[a4paper,margin=1in]{geometry}
\usepackage{amsmath,amssymb,amsthm,mathtools}
\usepackage{booktabs}
\usepackage{array}
\usepackage{xcolor}
\usepackage{microtype}
\usepackage{enumitem}
\usepackage{listings}
\usepackage[hidelinks]{hyperref}

\setlength{\parindent}{0pt}
\setlength{\parskip}{0.55em}
\setlength{\emergencystretch}{3em}

\newtheorem{lemma}{Lemma}[section]
\newtheorem{theorem}[lemma]{Theorem}
\newtheorem{proposition}[lemma]{Proposition}
\newtheorem{invariant}[lemma]{Invariant}
\theoremstyle{definition}
\newtheorem{definition}[lemma]{Definition}
\newtheorem{example}[lemma]{Example}
\theoremstyle{remark}
\newtheorem{remark}[lemma]{Remark}

\lstset{basicstyle=\ttfamily\small,columns=fullflexible,keepspaces=true,
  frame=single,framerule=0.3pt,xleftmargin=1em,xrightmargin=1em}

\newcommand{\Tid}{\mathsf{Tid}}
\newcommand{\Val}{\mathsf{Val}}
\newcommand{\dom}{\mathop{\mathrm{dom}}}
\newcommand{\Ready}{\mathsf{Ready}}
\newcommand{\Atomic}{\mathsf{Atomic}}
\newcommand{\Idle}{\mathsf{Idle}}
\newcommand{\Pend}{\mathsf{Pending}}
\newcommand{\Enq}{\mathsf{enq}}
\newcommand{\Deq}{\mathsf{deq}}
\newcommand{\Inc}{\mathsf{inc}}
\newcommand{\Read}{\mathsf{read}}
\newcommand{\Swap}{\mathsf{swap}}
\newcommand{\Store}{\mathsf{store}}
\newcommand{\ok}{\mathsf{ok}}
\newcommand{\nul}{\mathsf{null}}
\newcommand{\live}{\mathsf{live}}
\newcommand{\pending}{\mathsf{pending}}
\newcommand{\complete}{\mathsf{complete}}
\newcommand{\minimal}{\mathsf{minimal}}
\newcommand{\omitted}{\mathsf{omitted}}
\newcommand{\topo}{\mathsf{topo}}
\newcommand{\vals}{\mathsf{vals}}
\newcommand{\Scan}{\mathsf{Scan}}
\newcommand{\poss}{\Delta}
\newcommand{\conc}{\sigma}
\newcommand{\abst}{\rho}
\newcommand{\tokens}{\pi}
\newcommand{\entails}{\mathrel{\vDash}}
\newcommand{\code}[1]{\texttt{\detokenize{#1}}}
\newcommand{\Counter}{\mathsf{Counter}}
\newcommand{\SwapArray}{\mathsf{SwapArray}}
\newcommand{\POQueue}{\mathsf{POQueue}}
\newcommand{\Queue}{\mathsf{Queue}}
\newcommand{\HWQ}{\mathsf{HWQueue}}
\newcommand{\QLin}{\mathsf{QueueLin}}
\newcommand{\refines}{\rightsquigarrow}
\newcommand{\step}[1]{\xrightarrow{\;#1\;}}
\newcommand{\tok}[1]{\circ\,#1}
\newcommand{\done}[2]{\bullet\,(#1,#2)}
\newcommand{\lini}[1]{\star\,#1}
\newcommand{\ms}[1]{\mathsf{#1}}

\title{A Compositional Linearizability Proof of the\\Herlihy--Wing Queue in RGSimLin}
\author{Proof write-up for the \code{rbgs-rglin} development}
\date{October 2026}

\begin{document}
\maketitle

\begin{abstract}
The Herlihy--Wing queue is the standard example of an object whose
linearization points cannot be fixed in the code: the relative order of two
overlapping enqueues is decided by a later dequeue.  We show how the queue is
verified in the RGSimLin framework by the same decomposition used for the
timestamped stack.  An intermediate \emph{partial-order queue} specification
with interval-sequential enqueues isolates the future-dependent choice.  Below
it, the implementation is proved by a forward simulation with a single
abstract possibility and one loop invariant.  Above it, the partial-order
queue is linearized to the sequential FIFO queue by pure speculation over sets
of possibilities, with no reference to the code.  Vertical composition yields
linearizability of the implementation against the FIFO queue.
\end{abstract}

\tableofcontents

\section{The queue and why it is hard}

\subsection{The implementation}

The queue consists of an unbounded array \code{items} of cells, each holding
a value or \code{null}, and a counter \code{back}.  We use zero-based indices
and let \code{inc} return the \emph{old} counter value, which shifts the
original code by one but changes nothing else.

\begin{lstlisting}
enq(x):  i <- inc(back);           // allocate a slot
         store(items[i], x);       // fill it
         return ok

deq():   while true:
           r <- read(back);
           for i in 0 .. r-1:
             x <- swap(items[i], null);
             if x != null: return x
\end{lstlisting}

A dequeue scans the slots that existed when it read the counter and takes
the first non-empty one.  It never reports an empty queue; on an empty queue
it loops.  Section~\ref{sec:empty} explains why a single scan cannot be used
to return ``empty'' without breaking linearizability.

\subsection{Future-dependent linearization}

Let two enqueues $A$ and $B$ overlap: $A$ increments first and obtains slot
$0$, $B$ obtains slot $1$ but stores first.  A dequeue that scans slot $0$
while it is still empty takes $B$.  A dequeue that arrives after $A$ stored
takes $A$.  Which order the two enqueues take in the linearization is thus
decided by the first dequeue that touches them, possibly long after both
enqueues returned.  No single instruction of \code{enq} can serve as its
linearization point.

By contrast, a successful \code{swap} is the fixed linearization point of
its dequeue.  So there is exactly one source of future dependence, and it
concerns only the mutual order of overlapping enqueues.

\subsection{The decomposition}

Following the treatment of the TS stack in RGSimLin, we place a
\emph{partial-order queue} $\POQueue$ between the hardware objects and the
FIFO specification:
\[
\begin{array}{c}
\Queue \quad\text{(atomic FIFO queue)}\\[-2pt]
\uparrow\ M_{\QLin}\ \text{(forwarding; proof by speculation)}\\[-2pt]
\POQueue \quad\text{(interval-sequential enq, atomic deq)}\\[-2pt]
\uparrow\ M_{\HWQ}\ \text{(the code above; proof by forward simulation)}\\[-2pt]
\Counter \otimes \SwapArray
\end{array}
\]
$\POQueue$ records the enqueued elements as vertices of a DAG.  An enqueue is
\emph{interval sequential}: its invocation and its response are two separate
abstract events, and a vertex is ordered after exactly those vertices whose
enqueue had already responded when it was invoked.  Overlapping enqueues stay
unordered.  A dequeue is atomic and removes any minimal live vertex.

The implementation refines $\POQueue$ with fixed linearization points:
\code{inc} is the abstract enqueue invocation, \code{store} the abstract
enqueue response, and a successful \code{swap} the atomic dequeue.  The
difficulty of the queue then reappears, in a purely abstract form, as the
refinement of $\POQueue$ by the FIFO queue, where a set of possibilities keeps
every linear extension of the DAG alive until a dequeue selects one.

Unlike the TS stack, no trace theorem is needed: the stack's try-pop takes a
snapshot and is therefore not atomic with respect to its DAG specification,
whereas the queue's dequeue is.

\section{Preliminaries}

We use the notions of RGSimLin informally where the meaning is clear and
fix notation for the rest.

\paragraph{Events and specifications.}
Threads are ranged over by $\alpha, \beta \in \Tid$.  An observable event is
an invocation $\alpha{:}\,m$ or a response $\alpha{:}\,m.v$.  A specification
is a labelled transition system over such events.  An \emph{atomic} object is
given by a state set $S$ with transitions of the form
$\Idle(s) \step{\alpha:m} \Pend(s,\alpha,m) \step{\alpha:m.v} \Idle(s')$,
where the effect of the operation is taken at the response.  This is the
\code{VAE} construction of \code{AtomicLTS.v}.

\paragraph{Layers.}
A layer implementation $M : E \refines F$ provides, for each operation of
$F$, a program over the operations of $E$.  Its correctness statement is a
simulation establishing compositional linearizability, written
$M : E \refines F$ as well.  Simulations compose vertically
($M_1 \blacktriangleright M_2$) and horizontally ($M_1 \otimes M_2$), and the
identity $\mathrm{id}_E$ is a simulation.

\paragraph{Proof states and possibilities.}
A proof state consists of the concrete state $\conc$ of $E$ and an
\emph{abstract configuration} $\poss$: a nonempty set of possibilities
$(\abst, \tokens)$, where $\abst$ is a state of $F$ and $\tokens$ maps each
active thread to a \emph{token} recording how far its abstract operation has
progressed:
\[
\tok{m} \quad\text{invoked, not yet linearized};\qquad
\lini{m} \quad\text{abstract invocation taken, response pending};\qquad
\done{m}{v} \quad\text{linearized with result } v.
\]
For an atomic $F$ a possibility moves from $\tok{m}$ directly to
$\done{m}{v}$.  For the interval-sequential $\POQueue$ the two abstract
events of an enqueue are separate, and $\lini{\Enq(v)}$ marks the time in
between.  The simulation requires, at each concrete response, that
\emph{every} possibility carries the token $\done{m}{v}$ with the returned
$v$; at each concrete step the configuration may be updated by taking
abstract transitions in individual possibilities, by discarding possibilities
(filtering), and by adding new ones obtained from existing ones.  The
configuration must never become empty.

When a single possibility suffices we speak of a \emph{singleton} proof;
this is the facade of \code{SingletonPossibility.v}.  Otherwise the proof uses
the set logic of \code{RGILogicSet.v}.  Both supply rely--guarantee reasoning
with per-thread guarantees $G(\alpha)$ and relies
$R(\alpha) = \bigcup_{\beta \ne \alpha} G(\beta)$.

\section{The hardware objects}

\begin{definition}[Counter]
$\Counter$ is the atomic object with state $n \in \mathbb{N}$, initially
$0$, and operations
\[
\Inc : \quad n \mapsto n+1,\ \text{returns } n;
\qquad\qquad
\Read : \quad n \mapsto n,\ \text{returns } n.
\]
\end{definition}

\begin{definition}[Swap array]
$\SwapArray$ is the atomic object with state
$c : \mathbb{N} \to \Val \uplus \{\nul\}$, initially constantly $\nul$,
and operations
\[
\Swap(i, w) : \quad c \mapsto c[i \mapsto w],\ \text{returns } c(i);
\qquad\qquad
\Store(i, w) : \quad c \mapsto c[i \mapsto w],\ \text{returns } \ok.
\]
\end{definition}

Both objects are total and error-free.  Modelling the array as an infinite
function avoids any allocation reasoning: slot $i$ exists from the start and
is \emph{claimed} by the enqueue that receives $i$ from the counter.  The
underlay of the implementation is $E_0 = \Counter \otimes \SwapArray$.

\section{The partial-order queue \texorpdfstring{$\POQueue$}{POQueue}}
\label{sec:poqueue}

\subsection{State}

\begin{definition}[$\POQueue$ state]
A $\POQueue$ state is a tuple $s = (V, E, P, g, \mathit{nxt})$ with
\begin{itemize}[nosep]
\item $V : \mathbb{N} \rightharpoonup \Val$, the vertices and their values,
  including removed ones;
\item $E \subseteq \mathbb{N} \times \mathbb{N}$, the edges; $(n, m) \in E$
  reads ``$n$ is after $m$'';
\item $P : \Tid \rightharpoonup \mathbb{N}$, the pending enqueues: $P[\alpha]
  = n$ means thread $\alpha$ has invoked the enqueue of vertex $n$ and not
  yet responded;
\item $g \subseteq \mathbb{N}$, the removed (garbage) vertices;
\item $\mathit{nxt} \in \mathbb{N}$, the next vertex identifier.
\end{itemize}
The initial state is $(\emptyset, \emptyset, \emptyset, \emptyset, 0)$.
Derived predicates:
\begin{align*}
\pending(s, n) &\;:\Leftrightarrow\; \exists \alpha.\ P[\alpha] = n \\
\live(s, n) &\;:\Leftrightarrow\; n \in \dom V \wedge n \notin g \\
\complete(s, n) &\;:\Leftrightarrow\; n \in \dom V \wedge \neg\pending(s, n) \\
\minimal(s, n) &\;:\Leftrightarrow\; \live(s, n) \wedge
  \forall m.\ (n, m) \in E \Rightarrow \neg\live(s, m).
\end{align*}
\end{definition}

Vertices are identified by natural numbers handed out in order.  In the
implementation the identifier of a vertex will be its slot index, so this
choice removes a layer of indirection compared with the TS stack, whose
vertices are pairs of a thread and a location.

\subsection{Transitions}

As in the TryStack specification, a control state distinguishes
$\Ready(s)$, where enqueue events may happen, from $\Atomic(s, \alpha)$,
where only $\alpha$'s dequeue response is enabled.  The latter makes the
dequeue atomic: its invocation and response are adjacent in every trace.

\begin{definition}[$\POQueue$ transitions]\label{def:poq-steps}
\begin{align*}
\textsc{enq-inv}\quad&
\frac{P[\alpha] = \bot \qquad n = \mathit{nxt} \qquad
  E' = \{(n, m) \mid \complete(s, m)\}}
 {\Ready(V, E, P, g, \mathit{nxt}) \step{\alpha:\Enq(v)}
  \Ready(V[n \mapsto v],\ E \cup E',\ P[\alpha \mapsto n],\ g,\ \mathit{nxt}+1)}
\\[6pt]
\textsc{enq-ret}\quad&
\frac{P[\alpha] = n}
 {\Ready(V, E, P, g, \mathit{nxt}) \step{\alpha:\Enq(v).\ok}
  \Ready(V, E, P[\alpha \mapsto \bot], g, \mathit{nxt})}
\\[6pt]
\textsc{deq-inv}\quad&
\frac{}
 {\Ready(s) \step{\alpha:\Deq} \Atomic(s, \alpha)}
\\[6pt]
\textsc{deq-res}\quad&
\frac{\minimal(s, n) \qquad \complete(s, n) \qquad V(n) = v}
 {\Atomic((V, E, P, g, \mathit{nxt}), \alpha) \step{\alpha:\Deq.v}
  \Ready(V, E, P, g \cup \{n\}, \mathit{nxt})}
\end{align*}
\end{definition}

The edges created by \textsc{enq-inv} go to every complete vertex, garbage
included.  Edges into garbage are harmless, and making the rule independent
of $g$ keeps the well-formedness facts below simple (this is the
``garbage-independent edges'' decision of the TryStack development).

\subsection{Well-formedness}

\begin{lemma}[Well-formedness]\label{lem:wf}
Every reachable $\POQueue$ state satisfies
\begin{enumerate}[label=(W\arabic*),nosep]
\item $\dom V = \{ n \mid n < \mathit{nxt} \}$;
\item $(n, m) \in E \Rightarrow m < n$;
\item $\pending(s, m) \Rightarrow \forall n.\ (n, m) \notin E$;
\item $\pending(s, n) \Rightarrow n \notin g$;
\item $P[\alpha] = P[\beta] = n \Rightarrow \alpha = \beta$.
\end{enumerate}
\end{lemma}

\begin{proof}
By induction on the trace.  \textsc{enq-inv} adds the vertex
$\mathit{nxt}$ and edges from it to vertices in $\dom V$, which are smaller
by (W1); the new vertex is pending and receives no edge since $E'$ only has
it as a source; (W5) holds since $P[\alpha]$ was $\bot$ and the new identifier
is fresh.  \textsc{enq-ret} removes a pending entry and so can only weaken the
premises of (W3)--(W5).  \textsc{deq-res} adds to $g$ only a complete vertex,
preserving (W4), and leaves $V$, $E$, $P$ unchanged.
\end{proof}

A consequence of (W3) that drives the whole proof: a vertex can acquire
incoming edges only after its enqueue has responded.

\section{Layer 1: the implementation refines \texorpdfstring{$\POQueue$}{POQueue}}
\label{sec:layer1}

\subsection{Programs and linearization points}

$M_{\HWQ} : \Counter \otimes \SwapArray \refines \POQueue$ consists of the
two programs of Section~1, where \code{inc}, \code{read} address the counter
and \code{store}, \code{swap} the array.  The abstract events are taken as
follows, each at a single concrete step:

\begin{center}
\begin{tabular}{@{}lll@{}}
\toprule
concrete step & abstract events in the possibility & rule \\
\midrule
$\Inc$ returns $i$ & $\alpha{:}\Enq(x)$ with vertex $i$ & \textsc{enq-inv} \\
$\Store(i, x)$ & $\alpha{:}\Enq(x).\ok$ & \textsc{enq-ret} \\
$\Swap(i, \nul)$ returns $x \ne \nul$ & $\alpha{:}\Deq$ then $\alpha{:}\Deq.x$ & \textsc{deq-inv}, \textsc{deq-res} \\
$\Read$, $\Swap$ returning $\nul$ & none & \\
\bottomrule
\end{tabular}
\end{center}

Because every concrete step has at most one abstract counterpart and no
choice is ever revised, a singleton configuration suffices.  The proof is
carried out in the singleton facade exactly as the Timestamp and NodeMem
layers.

\subsection{The invariant}

Write $c$ for the counter value, $\mathit{cells}$ for the array, and
$s = (V, E, P, g, \mathit{nxt})$ for the abstract payload.

\begin{invariant}[Global invariant $I$]\label{inv:global}
\begin{enumerate}[label=(I\arabic*),nosep]
\item $s$ is well-formed (Lemma~\ref{lem:wf}) and $c = \mathit{nxt}$;
\item for all $i$ and $x \in \Val$:
  $\mathit{cells}(i) = x \iff \live(s, i) \wedge \complete(s, i) \wedge V(i) = x$;
\item the abstract control is $\Ready(s)$.
\end{enumerate}
\end{invariant}

(I2) can be read contrapositively: a slot is $\nul$ exactly when it has not
been handed out yet ($i \ge \mathit{nxt}$), its enqueue has not stored yet
(pending), or its element has already been dequeued (garbage).  (I3) holds at
every interference point because the two dequeue events are taken within one
concrete step.

\begin{definition}[Thread-local assertions]
For an enqueuing thread $\alpha$ with argument $x$:
\begin{align*}
\ms{EnqStart}(\alpha, x) &:\Leftrightarrow\; \tokens[\alpha] = \tok{\Enq(x)} \wedge P[\alpha] = \bot, \\
\ms{EnqSlot}(\alpha, x, i) &:\Leftrightarrow\; \tokens[\alpha] = \lini{\Enq(x)} \wedge P[\alpha] = i \wedge V(i) = x, \\
\ms{EnqDone}(\alpha, x) &:\Leftrightarrow\; \tokens[\alpha] = \done{\Enq(x)}{\ok} \wedge P[\alpha] = \bot.
\end{align*}
For a dequeuing thread $\alpha$ inside the scan, with range $r$ and next
index $i$:
\[
\Scan(\alpha, i, r) \;:\Leftrightarrow\;
\tokens[\alpha] = \tok{\Deq} \;\wedge\; i \le r \le \mathit{nxt} \;\wedge\;
\forall m < i.\ \live(s, m) \wedge \complete(s, m) \Rightarrow
  \forall n.\ (n, m) \in E \Rightarrow n \ge r.
\]
\end{definition}

The heart of the proof is $\Scan$.  It says: among the slots already
scanned, any element that is still present and complete can only be
\emph{preceded} (in the DAG) by vertices that were created after the
dequeuer read the counter, which are exactly the vertices the current round
will not look at.  The intuition is that if a complete, live $m$ below $i$
had a predecessor $n < r$, then $n$'s enqueue was invoked after $m$'s had
responded, hence after $m$'s store, hence $m$ was in its slot when the scan
passed it, and the scan would have taken it.  The state-based formulation
replaces this temporal argument by stability.

\subsection{Rely and guarantee}

$G(\alpha)$ allows the four kinds of steps in the table above, each paired
with its abstract events, plus the thread's own concrete invocation and
response, and requires $I$ to be preserved.  Every shared action of the
object is one of these, so $R(\alpha) = \bigcup_{\beta\ne\alpha} G(\beta)$
describes all interference.

\begin{lemma}[Stability]\label{lem:stable}
$I$, $\ms{EnqStart}$, $\ms{EnqSlot}$, $\ms{EnqDone}$ and $\Scan$ are stable
under $R(\alpha)$.
\end{lemma}

\begin{proof}
$I$ is preserved by each step of $G(\beta)$ by construction; we check the
interesting clause (I2).  An \textsc{enq-inv} by $\beta$ obtaining $i =
\mathit{nxt}$ leaves $\mathit{cells}(i) = \nul$ and makes $i$ pending, so
both sides of (I2) are false for $i$; $c$ and $\mathit{nxt}$ both advance.
A \textsc{enq-ret} by $\beta$ for slot $i$ sets $\mathit{cells}(i) = x =
V(i)$ and makes $i$ complete; by (W4) $i$ is live, so both sides become true.
A \textsc{deq-res} by $\beta$ at slot $i$ sets the cell to $\nul$ and adds
$i$ to $g$; both sides become false.

The enqueue assertions mention only $\tokens[\alpha]$, $P[\alpha]$ and
$V(i)$; other threads never modify these ($V$ is only extended at fresh
identifiers, and (W5) gives $\alpha$ sole ownership of $P[\alpha]$).

For $\Scan(\alpha, i, r)$ consider each step of $\beta \ne \alpha$:
\begin{itemize}[nosep]
\item \textsc{enq-inv} creating vertex $n = \mathit{nxt} \ge r$ with edges
  $(n, m)$.  Every new edge has source $n \ge r$.  Liveness and completeness
  of the old vertices are unchanged; the new vertex is $\ge r \ge i$, so the
  universally quantified clause does not apply to it.  Also $\mathit{nxt}$
  grows, keeping $r \le \mathit{nxt}$.
\item \textsc{enq-ret} completing a vertex $m$.  If $m < i$ and $m$ is live,
  the clause now demands that all edges into $m$ have source $\ge r$.  Before
  the step $m$ was pending, so by (W3) it has no incoming edges at all, and
  $E$ is unchanged.  The clause holds vacuously.
\item \textsc{deq-res} by $\beta$.  It removes a vertex from the live set,
  which only weakens the premise of the clause.
\item \textsc{read} and failed \textsc{swap} change nothing relevant. \qedhere
\end{itemize}
\end{proof}

\subsection{The method proofs}

\begin{lemma}[Enqueue]\label{lem:enq}
$\{\ms{EnqStart}(\alpha, x) \wedge I\}\ \code{enq}(x)\ \{\ms{EnqDone}(\alpha, x) \wedge I\}$
under $G(\alpha)$.
\end{lemma}

\begin{proof}
\emph{inc.}  The counter returns $i = c = \mathit{nxt}$ and advances.  Take
\textsc{enq-inv} in the possibility: its premise $P[\alpha] = \bot$ is
$\ms{EnqStart}$, and $n = \mathit{nxt} = i$.  Afterwards $V(i) = x$,
$P[\alpha] = i$, $\tokens[\alpha] = \lini{\Enq(x)}$, so $\ms{EnqSlot}(\alpha,
x, i)$.  (I1) holds as both counters advanced; (I2) at $i$ holds since the
cell is $\nul$ (it was never handed out before, by (I1) and (I2)) and $i$ is
pending.

\emph{store.}  The cell $i$ becomes $x$.  Take \textsc{enq-ret}: $P[\alpha]
= i$ by $\ms{EnqSlot}$.  Now $i$ is complete, live by (W4), and $V(i) = x$,
so (I2) holds at $i$.  The token becomes $\done{\Enq(x)}{\ok}$, giving
$\ms{EnqDone}$, which is what the concrete response requires.
\end{proof}

\begin{lemma}[Dequeue takes a minimal element]\label{lem:minimal}
Assume $I$ and $\Scan(\alpha, i, r)$ with $i < r$.  If $\Swap(i, \nul)$
returns $x \ne \nul$, then $\live(s, i)$, $\complete(s, i)$, $V(i) = x$ and
$\minimal(s, i)$.
\end{lemma}

\begin{proof}
The first three facts are (I2).  For minimality let $(i, m) \in E$ and
suppose $\live(s, m)$.  By (W2) $m < i$.  By (W3) $m$ is not pending, since it
has the incoming edge $(i, m)$; so $m$ is complete.  Now $\Scan$ applies to
$m < i$ and yields $i \ge r$, contradicting $i < r$.
\end{proof}

\begin{lemma}[Dequeue]\label{lem:deq}
$\{\tokens[\alpha] = \tok{\Deq} \wedge I\}\ \code{deq}()\
\{\exists x.\ \tokens[\alpha] = \done{\Deq}{x} \wedge I\}$ under $G(\alpha)$.
\end{lemma}

\begin{proof}
The outer loop invariant is $\tokens[\alpha] = \tok{\Deq} \wedge I$.

\emph{read.}  Returns $r = c = \mathit{nxt}$.  Establish $\Scan(\alpha, 0, r)$:
$0 \le r \le \mathit{nxt}$ and the universal clause is vacuous.

\emph{for loop.}  The loop invariant is $\Scan(\alpha, i, r)$ for the current
index $i$, with $i \le r$.  Consider the swap at $i < r$.
\begin{itemize}[nosep]
\item It returns $\nul$.  By (I2), $i \ge \mathit{nxt}$, or $i$ is
  pending, or $i \in g$.  The first is impossible since $i < r \le
  \mathit{nxt}$.  So $i$ is not both live and complete, and
  $\Scan(\alpha, i+1, r)$ follows from $\Scan(\alpha, i, r)$ with the clause
  for $m = i$ vacuous.  No abstract step.
\item It returns $x \ne \nul$.  By Lemma~\ref{lem:minimal} the premises of
  \textsc{deq-res} hold for $n = i$.  Take \textsc{deq-inv} and
  \textsc{deq-res} with value $x$ in the possibility, returning to
  $\Ready$ with $g' = g \cup \{i\}$.  The cell is now $\nul$ and $i \in g'$,
  so (I2) holds at $i$; nothing else changed.  The token is
  $\done{\Deq}{x}$ and the program returns $x$.
\end{itemize}
If the loop runs to $i = r$ without success, the invariant of the outer loop
still holds and the next round starts with a fresh $\Read$.  Linearizability
is a safety property, so non-termination on an empty queue is permitted.
\end{proof}

\begin{theorem}[Layer 1]\label{thm:layer1}
$M_{\HWQ} : \Counter \otimes \SwapArray \refines \POQueue$.
\end{theorem}

\begin{proof}
By the singleton rely--guarantee simulation theorem of RGSimLin, from
Lemmas~\ref{lem:stable}, \ref{lem:enq}, \ref{lem:deq}, the compatibility of
$G(\alpha)$ with $R(\beta)$ for $\alpha \ne \beta$ (each $G$ is included in
the other's $R$ by definition), and the initial state, in which the counter
is $0 = \mathit{nxt}$ and all cells are $\nul$ with $\dom V = \emptyset$.
\end{proof}

\section{Layer 2: \texorpdfstring{$\POQueue$}{POQueue} linearizes to the FIFO queue}
\label{sec:layer2}

\subsection{Specification and implementation}

\begin{definition}[FIFO queue]
$\Queue$ is the atomic object with state $q \in \Val^{*}$, initially
$\varepsilon$, and operations
\[
\Enq(v) : \quad q \mapsto q \cdot v,\ \text{returns } \ok;
\qquad\qquad
\Deq : \quad v \cdot q' \mapsto q',\ \text{returns } v.
\]
$\Deq$ has no response transition on the empty queue; a pending dequeue
on an empty queue simply never responds.
\end{definition}

$M_{\QLin} : \POQueue \refines \Queue$ forwards each operation:
\code{enq(v)} calls $\POQueue.\Enq(v)$ and \code{deq()} calls
$\POQueue.\Deq$.  All the content of this layer is in the proof, exactly as
for the TSStack layer over TryStack.

\subsection{Orders and justified possibilities}

The concrete state is now a $\POQueue$ control state; let $s$ be its payload.
The abstract configuration $\poss$ is a set of possibilities $(\abst,
\tokens)$ with $\abst$ a $\Queue$ state and $\tokens$ a token map over
$\Queue$ operations.

\begin{definition}[Omitted and linearized vertices]
A pending vertex is \emph{omitted} in $\tokens$ if its thread has not yet
linearized the enqueue:
\[
\omitted(s, \tokens, n) \;:\Leftrightarrow\;
\exists \alpha, v.\ P[\alpha] = n \wedge \tokens[\alpha] = \tok{\Enq(v)}.
\]
The \emph{linearized set} of a possibility is
$X(s, \tokens) = \{ n \mid \live(s, n) \wedge \neg\omitted(s, \tokens, n) \}$.
\end{definition}

Every complete live vertex is in $X$.  A pending live vertex is in $X$ iff
the possibility has already placed its enqueue (then its token is
$\done{\Enq(v)}{\ok}$).

\begin{definition}[Topological orders]
A list $N$ of distinct vertices is a \emph{topological order} of a set
$Y$, written $\topo(s, Y, N)$, if it enumerates exactly $Y$ and, whenever
$(n, m) \in E$ with $n, m \in Y$, $m$ occurs before $n$ in $N$.  Write
$\vals(s, N)$ for the list $V(n_1) \cdots V(n_k)$ when $N = n_1 \cdots n_k$.
\end{definition}

\begin{definition}[Justified possibility]
$(\abst, \tokens)$ is \emph{justified} by $N$ in $s$ if $\abst =
\Idle(\vals(s, N))$ and $\topo(s, X(s, \tokens), N)$.
\end{definition}

\subsection{The invariant}

\begin{invariant}[Layer 2 invariant $J$]\label{inv:layer2}
\begin{enumerate}[label=(J\arabic*),nosep]
\item \emph{Completeness.}  For every $Y$ with
  $\{n \mid \live(s,n) \wedge \complete(s,n)\} \subseteq Y \subseteq
  \{n \mid \live(s,n)\}$ and every $N$ with $\topo(s, Y, N)$, there is
  $(\abst, \tokens) \in \poss$ justified by $N$ with $X(s, \tokens) = Y$,
  and whose tokens agree with the \emph{phase} of every thread as described
  in (J2).
\item \emph{Token soundness.}  For every $(\abst, \tokens) \in \poss$ and
  thread $\alpha$:
  if $\alpha$ is between its concrete \code{enq}$(v)$ invocation and its
  $\POQueue$ \textsc{enq-ret}, then $\tokens[\alpha] \in \{\tok{\Enq(v)},
  \done{\Enq(v)}{\ok}\}$, and if $P[\alpha] = n$ then $V(n) = v$;
  if $\alpha$ is after its \textsc{enq-ret} and before its concrete response,
  then $\tokens[\alpha] = \done{\Enq(v)}{\ok}$;
  if $\alpha$ is between its concrete \code{deq} invocation and its
  $\POQueue$ dequeue, then $\tokens[\alpha] = \tok{\Deq}$;
  if $\alpha$ is after a $\POQueue$ dequeue returning $v$ and before its
  concrete response, then $\tokens[\alpha] = \done{\Deq}{v}$.
\item $s$ is well-formed.
\end{enumerate}
\end{invariant}

Completeness is the paper's $I_{\mathrm{graph}}$ adapted to the queue:
whatever subset of the pending enqueues one decides to count as already
linearized, and whatever linear extension of the DAG one picks for the
resulting elements, some possibility realizes it.  As in the TSStack proof,
no soundness property of $\abst$ is needed; it suffices that tokens are
right when a concrete response is taken, which is (J2), and that abstract
steps taken in a possibility are legal, which the updates below guarantee by
construction.

The phase information in (J2) is tracked by a ghost phase map, as in the
TryStack proof; its maintenance is routine and omitted.

\subsection{Possibility updates}

Four concrete events change the shared state; the client-facing invocations
and responses of $M_{\QLin}$ only move tokens.

\paragraph{Client invocations.}
A concrete \code{enq}$(v)$ invocation by $\alpha$ sets
$\tokens[\alpha] := \tok{\Enq(v)}$ in every possibility; a concrete
\code{deq} invocation sets $\tokens[\alpha] := \tok{\Deq}$.  (J1) is
unaffected because $P[\alpha] = \bot$, so $\alpha$ owns no vertex.

\begin{lemma}[\textsc{enq-inv}]\label{lem:l2-enqinv}
Let $\alpha$ take $\POQueue$ \textsc{enq-inv} with vertex $n$ and value
$v$, producing $s'$ from $s$.  Update $\poss$ to $\poss'$ as follows.  For
each $(\abst, \tokens) \in \poss$ justified by $N$ with $\abst =
\Idle(\vals(s,N))$:
\begin{enumerate}[label=(\alph*),nosep]
\item keep $(\abst, \tokens)$; $n$ is omitted in it;
\item for every list $M$ of distinct vertices omitted in $\tokens$
  (pending vertices of threads $\beta_1, \dots, \beta_k$ with tokens
  $\tok{\Enq(v_j)}$), add the possibility obtained by taking the atomic
  $\Queue$ enqueues of $v, v_1, \dots, v_k$ in this order in $\abst$ and
  setting $\tokens[\alpha] := \done{\Enq(v)}{\ok}$,
  $\tokens[\beta_j] := \done{\Enq(v_j)}{\ok}$.
\end{enumerate}
Then $J$ holds for $(s', \poss')$.
\end{lemma}

\begin{proof}
Every added possibility arises from legal $\Queue$ transitions on threads
whose tokens were $\tok{\cdot}$, so (J2) is preserved; the new vertex's
thread has token $\tok{\Enq(v)}$ in (a) and $\done{\Enq(v)}{\ok}$ in (b),
both allowed with $V'(n) = v$.

For (J1), let $Y'$ and $N'$ be as required for $s'$.  The live complete
vertices of $s'$ are those of $s$ (the new vertex is pending).

If $n \notin Y'$: $Y'$ is a valid choice for $s$, and $N'$ is topological
in $s$ because $E' \subseteq E \cup \{(n, \cdot)\}$ and $n \notin Y'$.
The possibility from (a) is justified by $N'$ in $s'$ with $X(s', \tokens) =
Y'$, as $n$ is omitted.

If $n \in Y'$: write $N' = N_1 \cdot n \cdot M$.  Every complete live vertex
$m$ has the edge $(n, m) \in E'$, so it occurs before $n$: all of it is in
$N_1$, and every vertex of $M$ is pending in $s$.  Let $Y = $ the set of
$N_1$.  Then $Y$ is a valid choice for $s$ and $N_1$ is topological in $s$
(a sublist of a topological order, with no edge involving $n$).  By (J1)
for $s$ there is $(\abst, \tokens) \in \poss$ justified by $N_1$ with
$X(s, \tokens) = Y$; in it, every vertex of $M$ is omitted.  Branch (b) with
this $M$ yields a possibility with $\abst' = \Idle(\vals(s,N_1) \cdot v \cdot
\vals(s, M)) = \Idle(\vals(s', N'))$ and linearized set $Y \cup \{n\} \cup M
= Y'$.
\end{proof}

Branch (b) linearizes other threads' pending enqueues ahead of time.  This
is sound because in the atomic $\Queue$ a thread with token
$\done{\Enq(v)}{\ok}$ merely has its response outstanding, and the token will
be exactly what its eventual concrete response needs.

\begin{lemma}[\textsc{enq-ret}]\label{lem:l2-enqret}
Let $\alpha$ take $\POQueue$ \textsc{enq-ret} for vertex $n$, producing
$s'$.  Update $\poss' = \{(\abst, \tokens) \in \poss \mid \tokens[\alpha] =
\done{\Enq(v)}{\ok}\}$.  Then $\poss' \ne \emptyset$ and $J$ holds.
\end{lemma}

\begin{proof}
In $s'$, $n$ is complete, so every admissible $Y'$ contains $n$, and
$\live$, $E$ are unchanged.  Given $Y'$ and $N'$ topological, the same $Y'$
is admissible in $s$ and (J1) gives $(\abst,\tokens) \in \poss$ justified by
$N'$ with $X(s,\tokens) = Y'$; since $n \in Y'$ is pending in $s$ it is not
omitted, so $\tokens[\alpha] = \done{\Enq(v)}{\ok}$ and the possibility
survives the filter.  In $s'$ the vertex is complete, so $X(s', \tokens) =
Y'$ still.  Nonemptiness follows from any such $Y'$, for instance the live
complete set with any topological order, which exists because $E$ is acyclic
by (W2).  (J2) holds because the surviving tokens are the required ones.
\end{proof}

\begin{lemma}[Dequeue]\label{lem:l2-deq}
Let $\alpha$ take $\POQueue$ \textsc{deq-inv} and \textsc{deq-res} removing
vertex $n$ with $V(n) = v$, producing $s'$ with $g' = g \cup \{n\}$.  Update
$\poss'$ to consist of, for each $(\Idle(v \cdot q), \tokens) \in \poss$ with
head value $v$ justified by $n \cdot N$, the possibility obtained by taking
the atomic $\Queue$ dequeue: $(\Idle(q), \tokens[\alpha \mapsto
\done{\Deq}{v}])$.  Then $\poss' \ne \emptyset$ and $J$ holds.
\end{lemma}

\begin{proof}
(J2): $\tokens[\alpha]$ was $\tok{\Deq}$ by (J2) for $s$, so the abstract
step is legal, and afterwards the token is $\done{\Deq}{v}$ as required.

(J1): let $Y'$ and $N'$ be admissible for $s'$.  The live vertices of $s'$
are those of $s$ minus $n$, and $n$ is complete and live in $s$.  Put $Y = Y'
\cup \{n\}$, admissible in $s$.  The list $n \cdot N'$ is topological for
$Y$ in $s$: edges among $Y'$ are respected by $N'$; an edge $(m, n)$ with $m
\in Y'$ requires $n$ before $m$, which holds; an edge $(n, m)$ with $m \in
Y'$ would need $m$ before $n$, but $m$ is live and $n$ is minimal, so no
such edge exists.  By (J1) for $s$ there is a possibility justified by $n
\cdot N'$, i.e.\ $\abst = \Idle(v \cdot \vals(s, N'))$; it is selected, and
after the dequeue it is justified by $N'$ in $s'$ with linearized set $Y'$
(omittedness is unchanged since $n$ is complete).  Taking $Y' = $ the live
complete set of $s'$ shows $\poss' \ne \emptyset$.
\end{proof}

\paragraph{Client responses.}
When $\alpha$'s concrete \code{enq} returns, (J2) states that every
possibility has $\tokens[\alpha] = \done{\Enq(v)}{\ok}$, which is precisely
the side condition of the response rule; the token is cleared.  Likewise a
concrete \code{deq} returning $v$ finds $\done{\Deq}{v}$ everywhere.

\begin{lemma}[Stability of $J$]
$J$ is stable under the rely of every thread.
\end{lemma}

\begin{proof}
Each interfering step is one of the four updates above or a client
invocation or response, and each preserves $J$ by the preceding lemmas.
The thread-local facts used in the method proofs are the token constraints
of (J2) for the thread itself, which no other thread's update modifies:
branch (b) of Lemma~\ref{lem:l2-enqinv} changes other threads' tokens only
from $\tok{\Enq(v)}$ to $\done{\Enq(v)}{\ok}$, which (J2) permits in the
phase before \textsc{enq-ret}.
\end{proof}

\begin{theorem}[Layer 2]\label{thm:layer2}
$M_{\QLin} : \POQueue \refines \Queue$.
\end{theorem}

\begin{proof}
By the set-of-possibilities simulation theorem of RGSimLin from the lemmas
of this section.  Initially $s$ has no vertices and $\poss = \{(\Idle(\varepsilon),
\emptyset)\}$, which satisfies (J1) with $Y = \emptyset$, $N = \varepsilon$.
\end{proof}

\section{Composition}

\begin{theorem}[Herlihy--Wing queue]
$M_{\HWQ} \blacktriangleright M_{\QLin} : \Counter \otimes \SwapArray
\refines \Queue$.  Hence every client program running the Herlihy--Wing
queue over an atomic counter and an atomic swap array is linearizable with
respect to the FIFO queue.
\end{theorem}

\begin{proof}
Vertical composition of Theorems~\ref{thm:layer1} and~\ref{thm:layer2}.
If the counter is itself implemented, say by a CAS loop over a CAS register
$C$, with a simulation $M_{\Counter} : C \refines \Counter$, horizontal
composition gives $(M_{\Counter} \otimes \mathrm{id}_{\SwapArray})
\blacktriangleright M_{\HWQ} \blacktriangleright M_{\QLin}$ over
$C \otimes \SwapArray$.
\end{proof}

\section{Remarks}

\subsection{Returning ``empty'' is not linearizable}
\label{sec:empty}

Suppose \code{deq} returned $\mathit{empty}$ when a scan finds no element.
Consider: $\Enq(A)$ runs to completion (slot $0$).  Then $D_1$ invokes
\code{deq}, reads $r = 1$, and is preempted before its swap.  $\Enq(B)$
runs to completion (slot $1$).  $D_2$ invokes \code{deq}, reads $r = 2$,
takes $A$ from slot $0$, returns $A$.  $D_1$ resumes, finds slot $0$ empty,
and returns $\mathit{empty}$.  In any linearization $D_1$ comes after
$\Enq(A)$, and $B$ is enqueued and never dequeued, so the queue is nonempty
whenever $D_1$ could be placed.  The original algorithm therefore loops, and
$\POQueue$ has no empty response.  A linearizable empty test needs an
additional stability check, which would be a separate layer.

\subsection{Why dequeues may take pending vertices in the TS stack but not here}

The \textsc{deq-res} premise $\complete(s, n)$ is provable at layer 1 because
a cell is non-$\nul$ only after its enqueue's \code{store}, which is the
\textsc{enq-ret}.  Keeping the premise makes $\POQueue$ a stronger interface
and slightly simplifies Lemma~\ref{lem:l2-deq}, where $n \in Y'$ then holds
for every admissible $Y'$.  Dropping it would also be sound, with the
Lemma~\ref{lem:l2-deq} argument choosing $Y = Y' \cup \{n\}$ regardless.

\subsection{Comparison with the TS stack}

The decomposition is the queue analogue of
$\mathrm{TSStack} \blacktriangleleft \mathrm{TryStack}$: the DAG
specification, the interval-sequential enqueue rules, the completeness
invariant and the three possibility updates are those of the TSStack proof
with ``top'' replaced by ``minimal'' and the stack by a list with insertion at
the tail.  What disappears is the TryStack layer, whose trace theorem
justifies linearizing a snapshot-based try-pop.  The queue's dequeue is
atomic in $\POQueue$ already, so the whole future dependence is handled by
Lemma~\ref{lem:l2-enqinv}.

\subsection{Planned Coq artefacts}

\begin{center}
\begin{tabular}{@{}ll@{}}
\toprule
file & content \\
\midrule
\code{examples/HWQ/CounterSpec.v} & $\Counter$ \\
\code{examples/Common/SwapSpec.v} & $\SwapArray$ (nat-indexed, no allocation) \\
\code{examples/HWQ/POQueueSpec.v} & Section~\ref{sec:poqueue}, with Lemma~\ref{lem:wf} \\
\code{examples/HWQ/QueueSpec.v} & the FIFO $\Queue$ \\
\code{examples/HWQ/HWQueue.v}, \code{HWQueueProof.v} & Section~\ref{sec:layer1}, singleton facade \\
\code{examples/HWQ/QueueLin.v}, \code{QueueLinProof.v} & Section~\ref{sec:layer2}, set logic \\
\code{examples/HWQ/HWQueue_TopLevel.v} & the composed simulation \\
\bottomrule
\end{tabular}
\end{center}

\end{document}
